\documentclass[a4paper,10pt,english]{article}

\usepackage[margin=2.3cm]{geometry}
\usepackage[utf8]{inputenc}
\usepackage[T1]{fontenc}
\usepackage[english]{babel}
\usepackage{amsmath}
\usepackage{amssymb}
\usepackage{amsthm}
\usepackage[colorlinks,urlcolor=blue,citecolor=blue,linkcolor=blue]{hyperref}

\newtheorem{proposition}{Proposition}[section]

\newcommand{\E}{\mathbb E}
\newcommand{\Prob}{\mathbb P}
\newcommand{\norm}[1]{\left\lVert #1\right\rVert}
\newcommand{\R}{\mathbb R}
\newcommand{\1}{\mathbf 1}

\setlength{\emergencystretch}{3em}

\title{Specification Search with Estimated Weights and Outcomes}
\author{Th\'eo Voldoire and Nic Fishman}
\date{}

\begin{document}

\maketitle

\begin{abstract}
Specification search over inverse-probability-weighted and augmented
inverse-probability-weighted estimators is a search over fitted nuisance
models as well as reported targets.  A support-indexed estimating equation
gives a common expansion for both estimators.  Its linear term is the usual
influence function corrected for nuisance estimation, while its quadratic
term is a covariance-aware polynomial in the nuisance score and derivative
fields.  Logistic IPW generally retains the linear nuisance channel.  At the
true propensity and outcome regressions, AIPW is orthogonal, its influence
function is common to every correctly specified support, and searched
contrasts begin with mixed propensity--outcome errors.  This note develops
the deterministic expansion and the exact AIPW remainder that form the basis
for a uniform stochastic approximation.
\end{abstract}

\section{A support-indexed M-estimator}

Let $O_1,\ldots,O_n$ be independent observations with law $P$, and write
$P_n f=n^{-1}\sum_{i=1}^n f(O_i)$ and
$\mathbb G_n f=\sqrt n(P_n-P)f$.  Candidate covariates are indexed by
$[p]$, and the researcher may fit any support $S$ in
\begin{equation*}
        \mathcal S_s=\{S\subseteq[p]:|S|\le s\}.
\end{equation*}
For each $S$, let $\eta_S\in\R^{d_S}$ collect the nuisance parameters and
let $m_S(O;\eta_S)\in\R^{d_S}$ be their estimating function.  Its
population root $\eta_S^\star$ and empirical root $\hat\eta_S$ satisfy
\begin{equation}\label{eq:nuisance-roots}
        Pm_S(O;\eta_S^\star)=0,
        \qquad
        P_nm_S(O;\hat\eta_S)=0.
\end{equation}
The support-specific target is the mean of a scalar estimating map
$q_S(O;\eta_S)$,
\begin{equation}\label{eq:target-definition}
        \theta_S^\star=Pq_S(O;\eta_S^\star),
        \qquad
        \hat\theta_S=P_nq_S(O;\hat\eta_S).
\end{equation}
This form includes a joint Z-estimator because $q_S(O;\eta_S)-\theta$ is
the estimating equation for $\theta$.

Write a dot for differentiation with respect to $\eta_S$ and set
\begin{equation}\label{eq:population-derivatives}
        H_S=P\dot m_S(O;\eta_S^\star),
        \qquad
        a_S=P\dot q_S(O;\eta_S^\star).
\end{equation}
Here $H_S$ is a $d_S\times d_S$ matrix and $a_S$ is a $d_S$-vector.  The
following deterministic expansion isolates the only first-order object
needed for all supports.

\begin{proposition}[Uniform M-estimation expansion]
\label{prop:first-order}
Suppose every $H_S$ is invertible and, for a constant $\kappa$,
\begin{equation*}
        \max_{S\in\mathcal S_s}\norm{H_S^{-1}}_{\mathrm{op}}\le\kappa.
\end{equation*}
On an event on which all empirical roots in \eqref{eq:nuisance-roots}
exist, put
\begin{align*}
 r_n&=\max_{S\in\mathcal S_s}\norm{\hat\eta_S-\eta_S^\star},\\
 \epsilon_{m,n}
 &=\max_{S\in\mathcal S_s}\sup_{0\le t\le1}
 \norm{P_n\dot m_S\{O;\eta_S^\star
 +t(\hat\eta_S-\eta_S^\star)\}-H_S}_{\mathrm{op}},\\
 \epsilon_{q,n}
 &=\max_{S\in\mathcal S_s}\sup_{0\le t\le1}
 \norm{P_n\dot q_S\{O;\eta_S^\star
 +t(\hat\eta_S-\eta_S^\star)\}-a_S}.
\end{align*}
Then, simultaneously over $S\in\mathcal S_s$,
\begin{equation}\label{eq:nuisance-first-order}
 \hat\eta_S-\eta_S^\star
 =-\frac1{\sqrt n}H_S^{-1}\mathbb G_nm_S(O;\eta_S^\star)
 +R_{\eta,S},
 \qquad
 \norm{R_{\eta,S}}\le\kappa\epsilon_{m,n}r_n,
\end{equation}
and
\begin{equation}\label{eq:target-first-order}
 \hat\theta_S-\theta_S^\star
 =\frac1{\sqrt n}\mathbb G_n\varphi_S+R_{\theta,S},
 \qquad
 \varphi_S
 =q_S(O;\eta_S^\star)-\theta_S^\star
 -a_S^\top H_S^{-1}m_S(O;\eta_S^\star),
\end{equation}
where
\begin{equation}\label{eq:first-order-remainder}
 |R_{\theta,S}|
 \le
 \bigl(\max_{T\in\mathcal S_s}\norm{a_T}\,
       \kappa\epsilon_{m,n}+\epsilon_{q,n}\bigr)r_n.
\end{equation}
\end{proposition}

\begin{proof}
The fundamental theorem of calculus and the empirical equation give
\begin{equation*}
 0=\frac1{\sqrt n}\mathbb G_nm_S(O;\eta_S^\star)
   +H_S(\hat\eta_S-\eta_S^\star)
   +\int_0^1\{P_n\dot m_S(O;\eta_{S,t})-H_S\}
      (\hat\eta_S-\eta_S^\star)\,dt,
\end{equation*}
where $\eta_{S,t}=\eta_S^\star+t(\hat\eta_S-\eta_S^\star)$.
Multiplication by $H_S^{-1}$ proves \eqref{eq:nuisance-first-order} and
its remainder bound.  Applying the same identity to $q_S$ yields
\begin{equation*}
 \hat\theta_S-\theta_S^\star
 =\frac1{\sqrt n}\mathbb G_nq_S(O;\eta_S^\star)
  +a_S^\top(\hat\eta_S-\eta_S^\star)+R_{q,S},
 \qquad |R_{q,S}|\le\epsilon_{q,n}r_n.
\end{equation*}
Substitution of \eqref{eq:nuisance-first-order} proves
\eqref{eq:target-first-order}--\eqref{eq:first-order-remainder}.
\end{proof}

The population derivative $a_S$ determines whether fitted nuisances enter
the target at first order.  When $a_S=0$, Proposition~\ref{prop:first-order}
reduces the leading term to $\mathbb G_n
q_S(O;\eta_S^\star)/\sqrt n$.  If this function is also common across the
searched supports, its empirical fluctuation cancels from every searched
contrast.  The first support-dependent term is then quadratic.

\section{The quadratic Taylor field}

The quadratic term is the analogue of the covariance-aware field for
linear-regression search.  It can be written without choosing a particular
nuisance model.  Let $\ddot m_S[u,u]$ denote the vector obtained by applying
the Hessian of each component of $m_S$ to $(u,u)$, and use the same notation
for the scalar Hessian of $q_S$.  Define
\begin{equation}\label{eq:u-definition}
        u_S=H_S^{-1}\mathbb G_nm_S(O;\eta_S^\star).
\end{equation}

\begin{proposition}[Second-order expansion]
\label{prop:second-order}
Suppose the maps $m_S$ and $q_S$ are three times continuously
differentiable on neighborhoods of their population roots.  Let $c_n\ge1$
satisfy $c_n/\sqrt n=o(1)$.  Use a star on a function or derivative to
mean evaluation at $\eta_S^\star$.  Assume uniformly over
$S\in\mathcal S_s$ that the inverse Jacobians and population derivatives
through order three are bounded.  Assume also that the empirical third
derivatives are $O_{\Prob}(1)$ on the line segment from $\eta_S^\star$ to
$\hat\eta_S$ and that
\begin{equation}\label{eq:second-order-field-size}
 \norm{\mathbb G_nm_S^\star}
 +|\mathbb G_nq_S^\star|
 +\norm{\mathbb G_n\dot m_S^\star}_{\mathrm{op}}
 +\norm{\mathbb G_n\dot q_S^\star}
 +\norm{\mathbb G_n\ddot m_S^\star}
 +\norm{\mathbb G_n\ddot q_S^\star}
 =O_{\Prob}(c_n),
\end{equation}
and that $\max_S\norm{\hat\eta_S-\eta_S^\star}
=O_{\Prob}(c_n/\sqrt n)$.  Then, uniformly over $S$,
\begin{equation}\label{eq:nuisance-second-order}
\begin{split}
 \hat\eta_S-\eta_S^\star
 ={}&-\frac{u_S}{\sqrt n}
 +\frac1nH_S^{-1}\left\{
       \mathbb G_n\dot m_S^\star u_S
       -\frac12P\ddot m_S^\star[u_S,u_S]
       \right\}
 +O_{\Prob}\left(\frac{c_n^3}{n^{3/2}}\right),
\end{split}
\end{equation}
and
\begin{equation}\label{eq:target-second-order}
 \hat\theta_S-\theta_S^\star
 =\frac{L_S}{\sqrt n}+\frac{Q_S}{n}
 +O_{\Prob}\left(\frac{c_n^3}{n^{3/2}}\right),
\end{equation}
where
\begin{align}
 L_S={}&\mathbb G_nq_S^\star-a_S^\top u_S,
 \label{eq:linear-field}\\
 Q_S={}&a_S^\top H_S^{-1}\left\{
       \mathbb G_n\dot m_S^\star u_S
       -\frac12P\ddot m_S^\star[u_S,u_S]
       \right\}
       -(\mathbb G_n\dot q_S^\star)^\top u_S
       +\frac12P\ddot q_S^\star[u_S,u_S].
 \label{eq:quadratic-field}
\end{align}
\end{proposition}

\begin{proof}
Put $\Delta_S=\hat\eta_S-\eta_S^\star$.  Taylor expansion of the empirical
nuisance equation around $\eta_S^\star$ gives
\begin{equation*}
 0=\frac1{\sqrt n}\mathbb G_nm_S^\star
   +H_S\Delta_S
   +\frac1{\sqrt n}\mathbb G_n\dot m_S^\star\Delta_S
   +\frac12P\ddot m_S^\star[\Delta_S,\Delta_S]
   +O_{\Prob}\left(\frac{c_n^3}{n^{3/2}}\right),
\end{equation*}
uniformly in $S$.
The empirical Hessian fluctuation and the third-order Taylor term are of
the displayed remainder order under \eqref{eq:second-order-field-size}.
Define
\begin{equation*}
 v_S=H_S^{-1}\left\{
       \mathbb G_n\dot m_S^\star u_S
       -\frac12P\ddot m_S^\star[u_S,u_S]\right\}.
\end{equation*}
Substitution of $-u_S/\sqrt n+v_S/n$ into the preceding display cancels
its terms of orders $n^{-1/2}$ and $n^{-1}$.  The remaining terms are
$O_{\Prob}(c_n^3/n^{3/2})$ uniformly in $S$.  The empirical Jacobian on
the segment joining $\eta_S^\star$ and $\hat\eta_S$ is
$H_S+o_{\Prob}(1)$ uniformly in $S$, so the mean value identity and
boundedness of $H_S^{-1}$ give \eqref{eq:nuisance-second-order}.  A Taylor
expansion of the target yields
\begin{equation*}
 \hat\theta_S-\theta_S^\star
 =\frac1{\sqrt n}\mathbb G_nq_S^\star
  +a_S^\top\Delta_S
  +\frac1{\sqrt n}(\mathbb G_n\dot q_S^\star)^\top\Delta_S
  +\frac12P\ddot q_S^\star[\Delta_S,\Delta_S]
  +O_{\Prob}\left(\frac{c_n^3}{n^{3/2}}\right).
\end{equation*}
Substitution of \eqref{eq:nuisance-second-order} proves
\eqref{eq:target-second-order}--\eqref{eq:quadratic-field}.
\end{proof}

The expansion depends on the data only through the joint empirical field
\begin{equation}\label{eq:joint-empirical-field}
 \left\{\mathbb G_nq_S^\star,\,
 \mathbb G_nm_S^\star,\,
 \mathbb G_n\dot q_S^\star,\,
 \mathbb G_n\dot m_S^\star:S\in\mathcal S_s\right\}.
\end{equation}
Replacing this field by a centered Gaussian field with the same covariance
turns \eqref{eq:linear-field}--\eqref{eq:quadratic-field} into a matched
Gaussian polynomial.  A uniform stochastic approximation amounts to
coupling this polynomial to the empirical Taylor field while controlling
both the Taylor remainder and the replacement error at the scale of the
maximum.  Pointwise M-estimation theory supplies neither control.

When $a_S=0$, the quadratic field simplifies to
\begin{equation}\label{eq:orthogonal-quadratic-field}
 Q_S=-(\mathbb G_n\dot q_S^\star)^\top
 H_S^{-1}\mathbb G_nm_S^\star
 +\frac12P\ddot q_S^\star
 \left[H_S^{-1}\mathbb G_nm_S^\star,
       H_S^{-1}\mathbb G_nm_S^\star\right].
\end{equation}
The first term is a covariance-aware product of a target-derivative score
and a nuisance score.  The second records the curvature of the target map.
Orthogonality removes the remaining nuisance-curvature terms because they
are multiplied by $a_S$.

\section{Logistic propensity scores and IPW}

Let $A\in\{0,1\}$ be treatment, $Y=A Y(1)+(1-A)Y(0)$ the observed
outcome, and $X\in\R^p$ the candidate covariates.  Suppose treatment is
unconfounded given $X$ and the true propensity
$e_0(X)=\Prob(A=1\mid X)$ is bounded away from zero and one.  Write
\begin{equation*}
        \mu_a(X)=\E\{Y(a)\mid X\},
        \qquad
        \theta_0=\E\{\mu_1(X)-\mu_0(X)\}.
\end{equation*}
For a support $S$, let $Z_S=(1,X_S^\top)^\top$ and fit the logistic model
\begin{equation}\label{eq:logistic-propensity}
        e_S(X;\gamma)=\Lambda(Z_S^\top\gamma),
        \qquad
        \Lambda(t)=\frac{e^t}{1+e^t}.
\end{equation}
The population and empirical logistic scores are
\begin{equation}\label{eq:logistic-score}
 P\{Z_S(A-e_S(X;\gamma_S^\star))\}=0,
 \qquad
 P_n\{Z_S(A-e_S(X;\hat\gamma_S))\}=0.
\end{equation}
Define
\begin{equation}\label{eq:ipw-map}
 q_S^{\mathrm{IPW}}(O;\gamma)
 =\frac{AY}{e_S(X;\gamma)}
  -\frac{(1-A)Y}{1-e_S(X;\gamma)},
 \qquad
 \hat\theta_S^{\mathrm{IPW}}
 =P_nq_S^{\mathrm{IPW}}(O;\hat\gamma_S).
\end{equation}
This is the Horvitz--Thompson form; normalized weights are covered by
adding their two normalizing equations to the stacked estimating system.
The corresponding population target is
\begin{equation*}
        \theta_S^{\mathrm{IPW},\star}
        =Pq_S^{\mathrm{IPW}}(O;\gamma_S^\star).
\end{equation*}

Put $e_S^\star(X)=e_S(X;\gamma_S^\star)$ and
\begin{align}
 J_S&=P\{e_S^\star(X)(1-e_S^\star(X))Z_SZ_S^\top\},
 \label{eq:fisher-information}\\
 a_S^{\mathrm{IPW}}
 &=-P\left[Z_S\left\{
       \frac{AY(1-e_S^\star)}{e_S^\star}
       +\frac{(1-A)Ye_S^\star}{1-e_S^\star}
       \right\}\right].
 \label{eq:ipw-derivative}
\end{align}
The logistic nuisance Jacobian is $H_S=-J_S$, while
$a_S^{\mathrm{IPW}}$ is the population derivative of the IPW target.
Proposition~\ref{prop:first-order} therefore gives the uniform influence
representation
\begin{equation}\label{eq:ipw-influence}
\begin{split}
 \sqrt n(\hat\theta_S^{\mathrm{IPW}}
             -\theta_S^{\mathrm{IPW},\star})
 ={}&\mathbb G_n
 \{q_S^{\mathrm{IPW}}(O;\gamma_S^\star)
       -\theta_S^{\mathrm{IPW},\star}\}\\
 &+(a_S^{\mathrm{IPW}})^\top J_S^{-1}
       \mathbb G_n\{Z_S(A-e_S^\star)\}+o_{\Prob}(1),
\end{split}
\end{equation}
provided the remainders in Proposition~\ref{prop:first-order} are
$o_{\Prob}(n^{-1/2})$ uniformly over the searched supports.  If the
logistic model contains $e_0$, then
$\theta_S^{\mathrm{IPW},\star}=\theta_0$.  The second term in
\eqref{eq:ipw-influence} nevertheless remains unless the IPW target is
orthogonal to propensity estimation.

Under random assignment with probability $\pi$, every support contains
the true intercept-only propensity and
\begin{equation}\label{eq:randomized-information}
        J_S=\pi(1-\pi)P(Z_SZ_S^\top).
\end{equation}
Thus the inverse sparse covariance blocks in the linear-regression theory
become inverse sparse Fisher-information blocks here.  The covariance-aware
product in \eqref{eq:orthogonal-quadratic-field} is recovered when the
linear derivative $a_S^{\mathrm{IPW}}$ vanishes; with prognostic covariates,
the shifted product retains the signal-dominated linear channel in
\eqref{eq:ipw-influence}.

\section{AIPW and orthogonality}

For candidate functions $e,\bar\mu_0,\bar\mu_1$, define the augmented map
\begin{equation}\label{eq:aipw-map}
\begin{split}
 q^{\mathrm{AIPW}}(O;e,\bar\mu_0,\bar\mu_1)
 ={}&\bar\mu_1(X)-\bar\mu_0(X)
 +\frac{A}{e(X)}\{Y-\bar\mu_1(X)\}\\
 &-\frac{1-A}{1-e(X)}\{Y-\bar\mu_0(X)\}.
\end{split}
\end{equation}
The same support may index the propensity and outcome-regression models.
When the nuisances use different supports, $S$ denotes their tuple.  Their
estimating equations can be stacked into $m_S$, so
Propositions~\ref{prop:first-order} and \ref{prop:second-order} apply
without changing the generic object.

The AIPW population map has an exact mixed remainder that does not depend
on parametric models.

\begin{proposition}[Exact mixed remainder]
\label{prop:aipw-mixed-remainder}
Let $\bar e$ be bounded away from zero and one, and let
$\bar\mu_0,\bar\mu_1$ be arbitrary square-integrable functions.  Under
unconfoundedness,
\begin{equation}\label{eq:aipw-bias-identity}
\begin{split}
 &Pq^{\mathrm{AIPW}}(O;\bar e,\bar\mu_0,\bar\mu_1)-\theta_0\\
 &\quad=\E\left[(\bar e(X)-e_0(X))
 \left\{
 \frac{\bar\mu_1(X)-\mu_1(X)}{\bar e(X)}
 +\frac{\bar\mu_0(X)-\mu_0(X)}{1-\bar e(X)}
 \right\}\right].
\end{split}
\end{equation}
Consequently, the population map equals $\theta_0$ if either the
propensity is correct or both outcome regressions are correct.  At
$(e_0,\mu_0,\mu_1)$ its derivative in every nuisance direction is zero.
\end{proposition}

\begin{proof}
Conditional expectation of \eqref{eq:aipw-map} given $X$ replaces
$A\{Y-\bar\mu_1(X)\}$ by
$e_0(X)\{\mu_1(X)-\bar\mu_1(X)\}$ and the corresponding control term by
$(1-e_0(X))\{\mu_0(X)-\bar\mu_0(X)\}$.  Subtraction of
$\mu_1(X)-\mu_0(X)$ and collection of the two outcome-regression errors
gives \eqref{eq:aipw-bias-identity}.  The identity is bilinear at the
truth, which proves both conclusions.
\end{proof}

The parametric case fits directly into the same-sample expansion.  For
$a\in\{0,1\}$, let
\begin{equation}\label{eq:linear-outcome-models}
        \mu_{a,S}(X;\alpha_{a,S})=Z_S^\top\alpha_{a,S}
\end{equation}
and take
\begin{equation*}
 \eta_S=(\gamma_S^\top,\alpha_{0,S}^\top,\alpha_{1,S}^\top)^\top.
\end{equation*}
Write $q^{\mathrm{AIPW}}(O;\eta_S)$ for the augmented map obtained from
these three models and $\1(\cdot)$ for the indicator function.  The
stacked empirical estimator solves
\begin{equation}\label{eq:stacked-aipw-scores}
\begin{split}
 P_n\{Z_S(A-e_S(X;\hat\gamma_S))\}&=0,\\
 P_n\{\1(A=a)Z_S(Y-Z_S^\top\hat\alpha_{a,S})\}&=0,
 \qquad a\in\{0,1\}.
\end{split}
\end{equation}
The in-sample AIPW estimator is
\begin{equation}\label{eq:in-sample-aipw}
 \hat\theta_S^{\mathrm{AIPW}}
 =P_nq^{\mathrm{AIPW}}
 \{O;e_S(\,\cdot\,;\hat\gamma_S),
       \mu_{0,S}(\,\cdot\,;\hat\alpha_{0,S}),
       \mu_{1,S}(\,\cdot\,;\hat\alpha_{1,S})\}.
\end{equation}

The quadratic field is explicit when every searched support contains the
true nuisance functions.  Put
\begin{align}
 K_{1,S}&=P\{e_0(X)Z_SZ_S^\top\},
 &
 K_{0,S}&=P\{(1-e_0(X))Z_SZ_S^\top\},
 \label{eq:outcome-information}\\
 \xi_{\gamma,S}
 &=J_S^{-1}\mathbb G_n\{Z_S(A-e_0(X))\},
 &
 \xi_{a,S}
 &=K_{a,S}^{-1}\mathbb G_n
   [\1(A=a)Z_S\{Y-\mu_a(X)\}].
 \label{eq:nuisance-score-fields}
\end{align}
Let $\dot q_{\gamma,S}^\star,\dot q_{0,S}^\star,\dot q_{1,S}^\star$
denote the derivatives of the AIPW map with respect to
$\gamma,\alpha_0,\alpha_1$, evaluated at the true nuisance functions.

\begin{proposition}[In-sample AIPW field]
\label{prop:aipw-field}
Suppose every searched logistic and linear model contains
$e_0,\mu_0,\mu_1$, and suppose the conditions of
Proposition~\ref{prop:second-order} hold for the stacked equations in
\eqref{eq:stacked-aipw-scores}.  Then, uniformly over $S$,
\begin{equation}\label{eq:aipw-quadratic-expansion}
 \hat\theta_S^{\mathrm{AIPW}}-\theta_0
 =(P_n-P)q_0+\frac{Q_S^{\mathrm{AIPW}}}{n}
 +O_{\Prob}\left(\frac{c_n^3}{n^{3/2}}\right),
\end{equation}
where $q_0=q^{\mathrm{AIPW}}(O;e_0,\mu_0,\mu_1)$ and
\begin{equation}\label{eq:aipw-quadratic-field}
\begin{split}
 Q_S^{\mathrm{AIPW}}
 ={}&(\mathbb G_n\dot q_{\gamma,S}^\star)^\top\xi_{\gamma,S}
    +(\mathbb G_n\dot q_{0,S}^\star)^\top\xi_{0,S}
    +(\mathbb G_n\dot q_{1,S}^\star)^\top\xi_{1,S}\\
 &+\xi_{\gamma,S}^\top
   (K_{0,S}\xi_{1,S}+K_{1,S}\xi_{0,S}).
\end{split}
\end{equation}
\end{proposition}

\begin{proof}
The stacked population Jacobian is
$-\operatorname{diag}(J_S,K_{0,S},K_{1,S})$, so
\eqref{eq:nuisance-second-order} gives
\begin{equation*}
 \sqrt n
 \begin{pmatrix}
 \hat\gamma_S-\gamma_S^\star\\
 \hat\alpha_{0,S}-\alpha_{0,S}^\star\\
 \hat\alpha_{1,S}-\alpha_{1,S}^\star
 \end{pmatrix}
 =
 \begin{pmatrix}
 \xi_{\gamma,S}\\ \xi_{0,S}\\ \xi_{1,S}
 \end{pmatrix}
 +O_{\Prob}\left(\frac{c_n^2}{\sqrt n}\right).
\end{equation*}
At the truth, Proposition~\ref{prop:aipw-mixed-remainder} gives
$a_S=0$, and $q_S^\star=q_0$ for every support.  The empirical derivative
term in \eqref{eq:orthogonal-quadratic-field} gives the first line of
\eqref{eq:aipw-quadratic-field}.  Applying
\eqref{eq:aipw-bias-identity} to the first-order nuisance expansions gives
\begin{equation*}
 n\{Pq^{\mathrm{AIPW}}(O;\hat\eta_S)-\theta_0\}
 =\xi_{\gamma,S}^\top
  (K_{0,S}\xi_{1,S}+K_{1,S}\xi_{0,S})
  +O_{\Prob}\left(\frac{c_n^3}{\sqrt n}\right),
\end{equation*}
which is the remaining term in \eqref{eq:aipw-quadratic-field}.
\end{proof}

The common empirical mean in \eqref{eq:aipw-quadratic-expansion} cancels
from every searched contrast.  The full same-sample dependence remains in
the joint score field
$(\xi_{\gamma,S},\xi_{0,S},\xi_{1,S},
\mathbb G_n\dot q_{\gamma,S}^\star,
\mathbb G_n\dot q_{0,S}^\star,
\mathbb G_n\dot q_{1,S}^\star)$.  Its covariance is retained by the matched
Gaussian field.  If only one nuisance model is correct, double robustness
still fixes the population target, but a support-dependent first-order
field can remain through the misspecified nuisance.  Double robustness and
orthogonality are therefore distinct for specification search.
When supports have different pseudo-true nuisance functions, the general
expansion in \eqref{eq:target-second-order} retains their population
landscape and support-dependent linear field without further changes.

\section{Uniform stochastic approximation}

The generic expansion reduces the uniform theory to one probabilistic
statement.  Let $G$ be a centered Gaussian version of the joint field in
\eqref{eq:joint-empirical-field} with the same covariance, and let
$L_S(G)$ and $Q_S(G)$ denote the polynomial in
\eqref{eq:linear-field}--\eqref{eq:quadratic-field} evaluated at $G$.
For a deterministic scale $r_{n,s}>0$, the coupling statement has the form
\begin{equation}\label{eq:desired-approximation}
 \max_{S\in\mathcal S_s}\hat\theta_S
 =\max_{S\in\mathcal S_s}
\left\{\theta_S^\star+\frac{L_S(G)}{\sqrt n}
                         +\frac{Q_S(G)}n\right\}
 +o_{\Prob}(r_{n,s}),
\end{equation}
where $r_{n,s}$ must match the order of the searched Gaussian fluctuation.
The covariance of $G$ retains the
dependence among propensity scores, outcome scores, and their derivatives;
no independent-coordinate reduction is built into
\eqref{eq:desired-approximation}.

A proof of \eqref{eq:desired-approximation} needs uniform localization of
all sparse nuisance fits, sparse conditioning of their information
matrices, positivity, moment bounds for the score and derivative classes,
and an invariance theorem for the maximum of the resulting degree-two
field.  Logistic separation must be excluded on the same event that gives
localization.  These are the nonlinear counterparts of
sparse overlap, moment concentration, and covariance matching in the
linear-regression theory.

The expansion also identifies the computational object.  For IPW it is a
shifted covariance-aware product of the logistic treatment score and a
weighted outcome score, with $J_S^{-1}$ coupling their coordinates.  For
orthogonal AIPW it is a mixed propensity--outcome polynomial, together with
the empirical derivative field in \eqref{eq:aipw-quadratic-field}.  Any
top-$s$ reduction requires
additional separability of these fitted-score geometries; the matched
Gaussian polynomial is the target approximation without that reduction.

\end{document}
